\lecture{7}{Wed 24 Jan 2024 13:30}{}

Let $(X, \|\cdot \|)$ be a normed space. By Hahn-Banach, for all $x \in X$, there exists $f \in X^*$ such that $\|f\| = 1$, such that $f(x) = \|x\|$. 

Note: We have a natural pairing

\begin{align*}
	 X \times X^* &\longrightarrow \mathbb{C} \\
	 (x, f) &\longmapsto \left\langle x, f \right\rangle  = f(x)
.\end{align*}
And in this language, $j(x) = \left\langle x, * \right\rangle $.

\begin{proposition}
	\label{prop:DoubleDualEmbedding}
	
	$j: X \to X^{* *}$ is a linear isometric embedding.
\end{proposition}
\begin{proof}
	Linearity is clear. We want to verify $\|j(x)\| = \|x\|$. Unwrap the definitions,
	\begin{align*}
		\|j(x)\| &= \sup \{|j(x) f|: f \in X^*, \|f\| \le  1\} \\
		&= \sup_{\|f\| \le  1} \{\|f\| \|x\|\} \le \|x\|
	.\end{align*} 
	By corollary Corollary~\ref{cor:NormFromDual}, there exists $\|f\| \le 1$ such that $|f(x)| = \|x\|$. Thus $|j(x)| = \|x\|$.
\end{proof}

\begin{corollary}
	\label{cor:NormedSpaceBanachCompletion}
	
	Every normed space $(X, \|\cdot \|)$ has a Banach space completion.
\end{corollary}
\begin{proof}
	Consider $X \to j(X) \rightarrow \overline{j(X)}  \rightarrow X^{**}$. Now $\overline{j(X)}$ is Banach, and $X$ embeds densely into $\overline{j(X)} $.
\end{proof}

\begin{example}
	$j$ is not surjective in general. Consider the spaces
	\[
		X = C_0(\mathbb{N}) \quad X^* \cong \ell^1(\mathbb{N}) \quad X^{* *} \cong \ell^{\infty }(\mathbb{N})
	.\] 
\end{example}

\begin{definition}[Reflexive Space]
	\label{defn:ReflexiveSpace}
	If $j$ is surjective, i.e.
	\[
		X \cong X^{* *}
	,\] 
	we call $X$ \textit{reflexive}.
\end{definition}
We know finite-dimensional linear spaces are reflexive, from linear algebra.
\todo{ Work through specific examples of non-reflexive spaces.}


\begin{theorem}[Lp spaces are reflexive]
	\label{thm:LpSpacesAreReflexive}
	 $\ell^p(\mathbb{N})$ is reflexive, for  $1 < p < \infty $.
\end{theorem}
Todo.
\todo{Study the proof and write it here.}

\begin{lemma}[Minkowski's Lemma]
	\label{lem:Minkowski}
	Let $(X, \|\cdot \|)$ be a normed space, let $V \subset X$ be open convex such that $0 \in V$. 

	Define $q = q_V: X \to  [0,\infty )$, by $q(x) = \inf \{t > 0: x \in  tV\} $. Then
	\begin{enumerate}
		\item $V = \{x \in X: q(x) < 1\} $.
		\item $q$ is a sublinear function.
	\end{enumerate}
	
\end{lemma}
\begin{remark}
	If $V = U_X(0,1)$, then $q(x) = \|x\|$.
\end{remark}
\begin{proof}
	First prove $V \subset  \{x \in X: q(x) < 1\} $. If $x \in V$, then there exists $\delta > 0$ such that $( 1 + \delta) x \in V$. Then $x \in \frac{1}{1 + \delta}V$. It follows that $q(x) \le \frac{1}{1 + \delta} < 1$.

	On the other side, if $q(x) < 1$, then there exists $t$, $q(x) \le t < 1$, such that $x \in t V$. But $t V \subset V$ since $V$ is convex: we have $x / t \in V$, so
	\[
		x = (1 - t) \cdot 0 + t \left( x \cdot t \right)  \in V
	.\] 

	Now prove that $q$ is sublinear. 
	Easily we can see that $q(\lambda x) = \lambda q(x)$.

	Set $s = q(x) + \varepsilon$, $t = q(y) + \varepsilon$. Then $\frac{x}{s} \in V, \frac{y}{t} \in V$, since $q(x / s) = q(x) / s = q(x) / (q(x) + \varepsilon) < 1$. Now by convexity,

	\[
		\frac{x + y}{ s + t} = \frac{s}{s + t} \frac{x}{s} + \frac{t}{s + t} \frac{y}{t} \in V
	.\] 

	Thus, $q(x + y) \le  s + t = q(x) + q(y) + 2 \varepsilon$.

\end{proof}

